% ===== 第 1 章 c01c：Chan 2014，原文 pp.10--15 =====
% [原文 p.10]
从拓扑上看，底空间 $B$ 就是 $\R^n$ 中的一个上半空间，它容许一个同时带有奇点与边界的整仿射流形结构。边界 $\partial B \subset B$ 的原像是 $X$ 中的一个非环面光滑超曲面 $D$。判别轨迹 $\Gamma \subset B$ 是位于超平面 $H$ 中的一个实余维二的热带子簇，我们称 $H$ 为 $B$ 中的\emph{墙}。按定义，SYZ 纤维化底空间内的墙，就是那些在 $X$ 中界定 Maslov 指标为零的全纯圆盘的拉格朗日环面纤维的轨迹。它把底空间分成不同的腔室，拉格朗日环面纤维在各个腔室中的 Floer 理论行为各不相同。在 Gross 纤维化的情形，墙 $H \subset B$ 平行于边界超平面 $\partial B$，并把底空间分成两个腔室。

现在考虑 $X$ 中由 SYZ 纤维化 $\mu$ 的纤维所界定、且与超曲面 $D$ 交于一点（重数为 1）的全纯圆盘的（虚拟）计数（即 Maslov 指标为二的圆盘）。当一个点穿过墙从一个腔室移动到另一个腔室时，由相应拉格朗日纤维所界定的全纯圆盘的虚拟数目（或亏格 0 开 Gromov--Witten 不变量）发生跳跃，表现出\emph{墙穿越现象}。Auroux [8, 9] 以及 Chan、Lau 与 Leung [20] 应用 Fukaya、Oh、Ohta 与 Ono [54] 发展的精细机器对此作了分析。注意在此情形不存在散射现象，因为只有一面墙。通过对底空间中每个腔室应用 SYZ 对偶纤维化构造，再按墙穿越公式把所得的各块拼接起来，我们就得到\emph{瞬子修正}的或 \emph{SYZ 镜像族} $\check X$，它以 $X$ 的 K\"ahler 模空间为参数 [20, 3]。所得结果与物理论证的预言 [35, 89] 一致。

这一 SYZ 镜像构造非常精确，因为它精确地告诉我们：$\check X$ 上哪个复结构对应于 $X$ 上任一给定的辛结构------镜像 $\check X$ 的定义方程是一个显式表达式，完全用 $X$ 的 K\"ahler 参数与圆盘计数不变量写出。例如，$X = K_{\PP^2}$ 的 SYZ 镜像是\footnote{更准确地说，$K_{\PP^2}$ 的 SYZ 镜像是 Landau--Ginzburg 模型 $(\check X, W = u)$；见下一节。}
\begin{equation}
  \check X = \Bigl\{\, (u,v,z_1,z_2) \in \C^2 \times (\C^\times)^2 \ \Big|\
  uv = 1+\delta(q) + z_1 + z_2 + \frac{q}{z_1 z_2} \,\Bigr\},
  \tag{1.6.1}
\end{equation}
其中 $q$ 是 K\"ahler 参数，度量 $\PP^2$ 上 $K_{\PP^2}$ 的零截面内一条射影直线的辛面积；并且
\begin{equation}
  1+\delta(q) = \sum_{k=0}^{\infty} n_k q^k
  \tag{1.6.2}
\end{equation}
是亏格 0 开 Gromov--Witten 不变量的生成级数。

此外，SYZ 构造自然地定义\emph{SYZ 映射}，它是从 $X$ 的 K\"ahler 模空间到 $\check X$ 的复模空间的一个映射。正如 Gross 与 Siebert 所猜测的，SYZ 镜像族应当用典范坐标写出；换言之，SYZ 映射应当给出镜像映射的逆。这一猜想对环面 Calabi--Yau 曲面与三维流形的一些证据在 [20, 108] 中给出，而 Chan、Lau 与 Tseng [21] 证明了当 $X$ 是紧环面 Fano 流形上典范线丛的全空间时的猜想。

% [原文 p.11]
最近，通过应用 orbifold 技巧，这一猜想对所有环面 Calabi--Yau 流形都在 [19] 中被证明了。

证明这些结果的主要挑战，是计算由 Fukaya、Oh、Ohta 与 Ono [55] 定义的开 Gromov--Witten 不变量。由于全纯圆盘的模空间通常是高度受阻的，这些不变量一般极难计算。目前可用的技巧很少（例如开/闭相等性、环面镜像定理、退化技巧等）。例如，(1.6.2) 中的不变量可以计算出来：
\[
  n_k = 1,\,-2,\,5,\,-32,\,286,\,-3038,\,35870,\,\dots
  \qquad (k=0,1,2,3,4,5,6,\dots),
\]
这与 [63] 中的周期计算一致。

我们应当提到，SYZ 构造也可以反向进行 [3]（另见 [18, 第 5 节]）。例如，从二次曲线丛 (1.6.1) 出发，可以用与 [62, 69] 类似的方法构造出 SYZ 纤维化。虽然在此情形判别轨迹具有实余维一，但人们仍可构造 SYZ 镜像，并且正如所期望的，这又给出环面 Calabi--Yau 三维流形 $K_{\PP^2}$\footnote{更准确地说，(1.6.1) 的 SYZ 镜像是 $K_{\PP^2}$ 中一个光滑超曲面的补集。}。

然而，在环面情形之外，尚不清楚如何以如此显式的方式实施 SYZ 构造。一个主要问题是开 Gromov--Witten 不变量的良定义性。只在少数几个非环面情形（见 Liu [121] 与 Solomon [151]）我们有良定义的开 Gromov--Witten 不变量理论\footnote{不过，Fukaya [51, 52] 最近有关于为紧 Calabi--Yau 三维流形定义圆盘计数不变量的工作。}。

\section{非 Calabi--Yau 情形下的 SYZ}
镜像对称被发现后不久，就被推广到非 Calabi--Yau 情形，特别是通过 Batyrev [10]、Givental [59, 60, 61]、Kontsevich [104]、Hori--Vafa [90] 以及许多其他人的工作推广到 Fano 流形。与 Calabi--Yau 情形不同，镜像不再由一个流形给出；取而代之，它被预言为一对 $(\check X,W)$，其中 $\check X$ 是非紧 K\"ahler 流形，$W : \check X \to \C$ 是全纯函数。在物理文献中，这样的对 $(\check X,W)$ 称为 Landau--Ginzburg 模型，$W$ 称为该模型的\emph{超势} [160, 168]。

自然会问：SYZ 方案在这一情形是否依然有效？Auroux [8] 是第一个考虑这一问题的人，事实上他把 SYZ 方案推广到了一个更为一般的情形。也就是说，他考虑由紧 K\"ahler 流形 $X$ 与一个有效反典范除子 $D$ 组成的对 $(X,D)$。$D$ 的定义截面给出 $X \setminus D$ 上沿除子 $D$ 具有单极点的全纯体积形式，因此在补集 $X \setminus D$ 上谈论特殊拉格朗日环面纤维化是有意义的。假设给定这样一个纤维化 $\mu : X \setminus D \to B$，那么我们可以尝试通过 T 对偶（即考虑对 $(L,\nabla)$ 的模空间，其中 $L$ 是 $\mu$ 的一个纤维，$\nabla$ 是 $L$ 上的平坦 $U(1)$-联络），并用瞬子修正来修改它，从而得到 SYZ 镜像 $\check X$。此外，超势 $W$ 将自然地作为 Fukaya--Oh--Ohta--Ono 的障碍链 $m_0$ 的镜像对象出现。

% [原文 p.12]
当 $X$ 是紧环面 K\"ahler 流形时，$D$ 的一个典范选择是所有环面素除子的并。此外，矩映射在 $X$ 上给出一个方便的拉格朗日环面纤维化，它有一个良好的性质：在开稠密的环面轨道 $X \setminus D$ 上限制为一个环面丛。在此情形，SYZ 镜像流形 $\check X$ 简单地由代数环面 $(\C^\times)^n$ 给出，因为我们有一个环面丛，且在镜像流形的构造中没有瞬子修正。所有本质信息都编码在超势 $W$ 中。在 Auroux 的工作之前，是 Cho 与 Oh [36, 38] 首先注意到 $W$ 可以用圆盘计数不变量（或开 Gromov--Witten 不变量）来表达。通过对 $X$ 中由矩映射纤维所界定的所有全纯圆盘进行分类，他们在 $X$ 为 Fano 的情形得到了 $W$ 的显式公式；这与 Hori 与 Vafa [90] 用物理论证早先得到的预言一致。后来，Fukaya、Oh、Ohta 与 Ono [55, 56, 53] 关于拉格朗日 Floer 理论与环面流形镜像对称的工作极大地推广了这一结果。

在 [22] 中，环面 Fano 流形的镜像对称被用作试验场，以检验 Fourier--Mukai 型变换（也就是我们所说的 SYZ 变换）在研究镜像对称几何时有多大用处。对于环面 Fano 流形 $X$，我们考虑开稠密的环面轨道 $X_0 := X \setminus D \subset X$，它也是矩映射的拉格朗日环面纤维的并。在辛几何上，我们可以写出 $X_0 = T^*B_0/\Lambda^\vee$，其中 $B$ 是矩多面体，$B_0$ 表示其内部。于是 SYZ 镜像是 $\check X := TB_0/\Lambda$，它是 $(\C^\times)^n$ 中的一个有界区域。为得到超势 $W$，我们考虑空间
\[
  \tilde X := X_0 \times \Lambda \subset \mathcal{L}X,
\]
即 $X$ 中逐纤维的测地/仿射环路构成的空间。在 $\tilde X$ 上，我们有一个瞬子修正的辛结构 $\tilde\omega = \omega + \Phi$，其中 $\Phi$ 是亏格 0 开 Gromov--Witten 不变量的生成函数，这些不变量（虚拟地）计数由矩映射纤维界定的全纯圆盘。

于是，通过把半平坦 SYZ 变换 $\mathcal{F}^{\text{semi-flat}}$ 与逐纤维 Fourier 级数结合起来，可以构造出显式的 SYZ 变换 $\mathcal{F}$；并已证明 $\mathcal{F}$ 把 $X$ 上修正后的辛结构 $\tilde\omega$ 恰好变为镜像 Landau--Ginzburg 模型 $(\check X,W)$ 的全纯体积形式 $e^{W}\check\Omega$，其中 $W$ 由对 $\Phi$ 取逐纤维 Fourier 变换得到。此外，$\mathcal{F}$ 诱导出 $X$ 的（小）量子上同调环 $QH^*(X)$ 与 $W$ 的 Jacobi 环 $\operatorname{Jac}(W)$ 之间的同构。证明的办法是取热带极限，并观察到：其全纯对应物贡献于量子积的热带曲线，可以由热带圆盘粘接得到（更多细节见 [22]）。这一观察后来被 Gross [71] 推广，用于他借助热带几何对 $\PP^2$ 的大量子上同调镜像对称的研究。

至于一般型的流形，目前沿 SYZ 观点研究其镜像对称有两条主要途径。一条是 Abouzaid、Auroux 与 Katzarkov 的工作：他们考虑环面簇 $V$ 中的超曲面 $H$，并构造一个 Landau--Ginzburg 模型，它是 $V \times \C$ 沿 $H \times \{0\}$ 的爆破的 SYZ 镜像。特别地，当 $H$ 是 $V = \PP^1 \times \PP^1$ 中双次数 $(3,2)$ 多项式的零集时，他们的构造给出亏格 2 Riemann 曲面的一个镜像，这与 Katzarkov 早先的提法 [100, 98, 147] 一致。

另一条途径与 Gross--Siebert 纲领更一致，是 Gross、Katzarkov 与 Ruddat [74] 的工作。他们提出，一般型簇的镜像是一个可约簇，并带有某个消没轮层。

% [原文 p.13]
大概说来，这一途径产生的镜像应当给出与 [3] 所产生者相同的数据。例如，该可约簇应当是 SYZ 镜像 Landau--Ginzburg 模型的超势的临界轨迹。但这两条途径之间的精确关系仍在研究中。

\section{SYZ 之外}
除了为镜像对称提供一个优美的几何解释之外，SYZ 猜想 [154] 还对数学的许多相关领域产生了持久影响。让我们简要描述这方面的几个应用例子。

\para{借助 SYZ 的同调镜像对称。}正如我们所见，SYZ 猜想建立在弦论中 D-膜的想法之上。回忆一下：B-膜（即 B-模型中的 D-膜）是复子簇上的凝聚层，而 A-膜（即 A-模型中的 D-膜）是配备平坦 $U(1)$ 联络的特殊拉格朗日子流形。因此，把 Kontsevich 的 HMS 猜想 [103]（它断言 Calabi--Yau 流形 $X$ 的 Fukaya 范畴等价于镜像 $\check X$ 上凝聚层的导出范畴）看作 $X$ 上 A-模型与 $\check X$ 上 B-模型之间同构的表现，是合理的。于是很自然地，人们期望可以利用 SYZ 方案、特别是 SYZ 变换，来构造实现 HMS 猜想所断言的范畴等价的函子。

例如，给定特殊拉格朗日环面纤维化 $\mu : X \to B$ 的一个拉格朗日截面，它与 $\mu$ 的纤维 $L$ 的交点确定对偶环面 $L^\vee$ 上的一个平坦 $U(1)$-联络。把这些平坦 $U(1)$-联络拼接起来，应当给出对偶纤维化全空间（即镜像 $\check X$）上的一个全纯线丛。这个简单的想法最先由 Gross [67, 68] 设想，继而被 Arinkin 与 Polishchuk [7] 以及 Leung、Yau 与 Zaslow [113] 用来构造 SYZ 变换，并应用于证明和理解半平坦 Calabi--Yau 情形下的 HMS 猜想。后来，同样的想法也被用于研究环面簇的 HMS 猜想 [1, 2, 43, 45, 44, 16, 24, 37]。

在一些更近期的工作 [18, 26, 25] 中，SYZ 变换被用来构造（在对象层面上的）几何 Fourier 型函子，它们实现了某些局部 Calabi--Yau 例子（如 $A_n$ 奇点的解消与光滑化锥形流形）的 HMS 范畴等价；在这些例子中会遇到带奇异纤维的 SYZ 纤维化，从而产生非平凡的量子修正。另一方面，K.-L. Chan、Leung 与 Ma 正在进行的工作 [28, 27] 表明，SYZ 变换也可以用来在态射层面上构造 HMS 等价，至少在半平坦情形如此。最终目标是：对任意给定的 SYZ 纤维化构造出典范的几何 Fourier 型函子，实现 HMS 猜想所断言的范畴等价，从而加深我们对 HMS 猜想之几何、乃至整个镜像对称的理解。

\para{Ricci 平坦度量与圆盘计数。}SYZ 论文 [154] 中的一个引人注目的观察是：镜像上的 Ricci 平坦度量可以分解为半平坦部分（此前已在 [65] 中显式写出）与瞬子修正部分之和；后者

% [原文 p.14]
应当来自原 Calabi--Yau 流形中以 SYZ 纤维化的拉格朗日环面纤维为边界的全纯圆盘的贡献。这给出了 Ricci 平坦度量的一种具体而定性的描述，而这类度量原本极难写出。

一般说来，这样的定性描述仍然极难获得，因为在紧 Calabi--Yau 流形上很难找到非平凡的 SYZ 纤维化例子，而开 Gromov--Witten 理论尚未被很好地理解。不过，Gaiotto、Moore 与 Neitzke 近期的开创性工作 [57, 58] 为超 K\"ahler 情形带来了新的曙光。他们提出了复可积系统全空间上的超 K\"ahler 度量（其最简单的例子重现了著名的 Ooguri--Vafa 度量 [139]）与 Kontsevich--Soibelman 的墙穿越公式之间的一种新关系（部分仍是猜想）。

为稍微详细地描述他们的提法，我们考虑一个复可积系统 $\psi : M \to B$，即 $M$ 是全纯辛的，且 $\psi$ 的纤维是复拉格朗日子流形。（更准确地说，Gaiotto、Moore 与 Neitzke 在 [57, 58] 中考察的是亚纯 Hitchin 系统，在此情形完全的 超 Kähler 度量最先由 Biquard 与 Boalch [11] 构造；见 [58, 第 4.1 节]。）他们利用了这样一个事实：任何 超 Kähler 度量都由其相伴的 扭量空间刻画；并尝试写出 $M$ 上一族 $C^\times$-全纯 Darboux 坐标，使其满足 Hitchin 等人定理 [88] 的假设。特别地，他们要求这些坐标满足某些墙穿越公式，这些公式描述坐标在所谓 BPS 射线两侧的不连续性；在 BPS 射线上，BPS 态的（虚拟）计数发生跳跃。

这些墙穿越公式与 Kontsevich 和 Soibelman [107] 以及 Gross 与 Siebert [79] 在构造 Calabi--Yau 流形环面退化时所使用的公式属于同一类型（另一方面，它们也与 motivic Donaldson--Thomas 理论 [97, 105] 中的墙穿越公式相同）。鉴于这一点与 SYZ 猜想（以及这一事实：通过 [11] 中的 超 Kähler 旋转，人们已在亚纯 Hitchin 系统上构造出大量非紧 SYZ 纤维化的例子，包括复维数二情形中的许多例子，如引力瞬子、对数 Calabi--Yau 曲面），自然地期望 Gaiotto、Moore 与 Neitzke 所考察的那些复可积系统上的 超 Kähler 度量可以用全纯圆盘来表达。

这一点在 [17] 中对最简单的例子------Ooguri--Vafa 度量------已经做到了。W. Lu 更近期的工作 [125, 126] 表明：一般地，[57, 58] 所研究的亚纯 Hitchin 系统上的 扭量空间与全纯 Darboux 坐标所产生的数据，与运行 Gross--Siebert 纲领 [79] 所需的数据相同；因此 超 Kähler 度量与热带圆盘计数之间必定存在某种（也许是隐含的）关系。在他博士论文 [120] 中，Y.-S. Lin 考察了椭圆 K3 曲面与热带圆盘计数不变量。他证明他的不变量满足与 [57, 58] 中出现的相同的墙穿越公式。这再次表明这些 K3 曲面上的 超 Kähler 度量与圆盘计数密切相关。Stoppa 及其合作者近期的文章 [152, 46] 也展示了 motivic Donaldson--Thomas 理论中的墙穿越公式与 Gaiotto、Moore 和 Neitzke 的构造之间的密切关系。

\para{SYZ 的其他应用。}我们还要提及 SYZ 构造最近两项出人意料的应用。

% [原文 p.15]（正文延伸至参考文献标题之前）
在其近期的合作项目 [72] 中，Gross、Hacking 与 Keel 通过把 [79] 中的构造推广到允许仿射流形具有更一般（即更坏）的奇点类型，构造了对数 Calabi--Yau 曲面（即满足 $Y$ 为非奇异射影有理曲面、$D \in |{-K_Y}|$ 是有理曲线之圈的对 $(Y,D)$）的镜像族。令人惊讶的是，他们的结果可以被用来证明 Looijenga [124] 一个三十年之久的猜想------关于尖点奇点的可光滑化性。

在甚至更近的一篇预印本 [73] 中，他们再次应用自己的构造，证明了对数 Calabi--Yau 曲面的 Torelli 定理；这一定理最初由 Friedman [49] 在 1984 年猜想。另一方面，他们的构造也紧密联系于 cluster 簇的理论，并且他们提出了 Fock--Goncharov 对偶基的一个巨大推广。关于这些激动人心的新结果与新进展的精彩阐述，我们请读者参阅 Gross 与 Siebert 的综述 [75]。

在另一个出人意料的方向上，SYZ 构造最近被应用于构造新的\emph{纽结不变量}。对 $S^3$ 中的纽结 $K$，其余法丛 $N^*K$ 典范地是余切丛 $T^*S^3$ 中的拉格朗日圈。在 [40] 中，Diaconescu、Shende 与 Vafa 在解消锥形流形 $X := \mathcal{O}_{\PP^1}(-1) \oplus \mathcal{O}_{\PP^1}(-1)$ 中构造了相应的拉格朗日圈 $L_K$；粗略地说，这是把余法丛 $N^*K$ 从零截面上抬起来，并让 $T^*S^3$ 经历锥形跃迁而得到的。他们的构造受物理中一个称为\emph{大 $N$ 对偶}的神秘现象启发。

在 [6] 中，Aganagic 与 Vafa 通过对 $(X,L_K)$ 应用推广的 SYZ 构造，定义了一个新的纽结不变量。粗略地说，他们的不变量是 $(X,L_K)$ 的开 Gromov--Witten 不变量的生成级数。结果表明，所得的不变量总是多项式，并且他们猜想它应当是纽结理论中经典 A-多项式的形变 [39]。此外，人们还提出了他们的不变量与 $K$ 的接触同调代数的增广 [135] 之间的一个有趣关系。关于这一关系的大量证据在一篇非常近期的预印本 [5] 中得到。

SYZ 猜想的这两项新应用，连同许多尚未到来的应用，为镜像对称以及数学与物理的其他分支开辟了新的方向\footnote{一个我们尚未提及的有趣故事，是 Gukov、Yau 与 Zaslow [85] 把 SYZ 图景应用于 $G_2$ 流形，其目的是解释 M-理论与杂化弦理论之间的对偶。}，并且它们都指向未来的进一步研究工作。
